\documentclass[a4paper, 12pt]{report}

%%%%%%%%%%%%
% Packages %
%%%%%%%%%%%%

\usepackage[english]{babel}
\usepackage[noheader]{sleek}
\usepackage{sleek-title}
\usepackage{sleek-listings}
\usepackage{booktabs}
\usepackage{longtable}

%%%%%%%%%%%%%%
% Title-page %
%%%%%%%%%%%%%%

\logo{./report/resources/pdf/Kcl-logo.pdf}
\faculty{Faculty of Natural, Mathematical, and Engineering Sciences}
\department{Department of Informatics}
\title{Developer's Manual}
\subtitle{StudyFlax: Time Management for Higher Education}
\author{\textbf{Team Flax-10}\\
Aishah Ahmed\\
Eric Almeida\\
Zerun (Harry) Chen\\
Cate Tay (Cate) Chew\\
Indixora Jozal\\
Kausar Khan\\
Hamed Latif\\
Dian Mahmood\\
Binuli Nevanma Markawitage
}
\date{\today}

%%%%%%%%%%%%
% Document %
%%%%%%%%%%%%

\begin{document}

\maketitle
\romantableofcontents

% ─────────────────────────────────────────────────────────────
\chapter{Overview}

StudyFlax is a web-based academic time management platform built for
students and teachers at King's College London. It allows students to
organise their schedules, track deadlines, and manage their academic
workload in one place. Teachers can create modules, set deadlines, and
assign work to their students, whilst the admin oversees and manages
the system as a whole.

The system consists of two software components:

\begin{itemize}
    \item \textbf{Backend} --- A Django REST API developed primarily
          with Python 3.12 in the Nix environment and backed
          by a PostgreSQL database.
    \item \textbf{Frontend} --- A React (Vite) single-page application.
\end{itemize}

The deployed application is available at:
\url{https://flax-10-production.up.railway.app/}

% ─────────────────────────────────────────────────────────────
\chapter{Prerequisites}

The following tools must be installed before setting up the development
environment. Installation of the operating system itself is assumed.

\begin{itemize}
    \item \textbf{Nix Package Manager} (v2.33.2 or later)\\
          \url{https://nixos.org/download/\#download-nix}

    \item \textbf{Git}\\
          \url{https://git-scm.com/downloads}

    \item \textbf{Python} (v3.12 recommended for manual local setup)\\
          \url{https://www.python.org/downloads/}

    \item \textbf{Node.js} (v20 recommended)\\
          \url{https://nodejs.org/}

    \item \textbf{PostgreSQL} (only required for manual setup without Nix)\\
          \url{https://www.postgresql.org/download/}
\end{itemize}

% ─────────────────────────────────────────────────────────────
\chapter{Cloning the Repository}

Clone the repository and navigate into it:

\begin{lstlisting}[language=bash]
git clone https://github.com/hoxtl/Flax-10.git
cd Flax-10
\end{lstlisting}

% ─────────────────────────────────────────────────────────────
\chapter{Setting Up the Development Environment with Nix}

This project uses a Nix flake to provide a reproducible development
environment. The Nix shell itself only provides the required tools and
prints the available entrypoints. Project setup and common tasks are
then run explicitly using the Nix packages defined in
\texttt{flake.nix}.

\begin{itemize}
    \item \texttt{nix develop} enters the shell and exposes Python,
          Node.js, Git, and PostgreSQL tooling.
    \item \texttt{nix run .\#init} creates the project virtual
          environment, installs Python and Node dependencies, prepares
          a project-local PostgreSQL instance, applies migrations, and
          seeds the database.
    \item \texttt{nix run .\#run} starts the backend and frontend
          development servers.
    \item \texttt{nix run .\#tests} runs backend and frontend automated
          tests and writes coverage reports to \texttt{./coverage/}.
    \item \texttt{nix run .\#seed} and \texttt{nix run .\#unseed}
          manage seeded data without reinitialising the whole project.
\end{itemize}

\section{Linux and macOS}

\begin{enumerate}
    \item Navigate to the root of the repository:
\begin{lstlisting}[language=bash]
cd /path/to/Flax-10
\end{lstlisting}

    \item Enter the Nix development shell:
\begin{lstlisting}[language=bash]
nix --extra-experimental-features 'nix-command flakes' develop
\end{lstlisting}

    \item \textbf{Optional --- enable experimental features permanently.}
          Add the following line to \texttt{/etc/nix/nix.conf} to avoid
          passing the flags on every invocation:
\begin{lstlisting}
experimental-features = nix-command flakes
\end{lstlisting}

    \item Initialise the repository:
\begin{lstlisting}[language=bash]
nix run .#init
\end{lstlisting}

    \item Start the development servers when required:
\begin{lstlisting}[language=bash]
nix run .#run
\end{lstlisting}

    \item Run the full automated test suite when required:
\begin{lstlisting}[language=bash]
nix run .#tests
\end{lstlisting}
\end{enumerate}

\section{Re-entering the Nix Shell in Future Sessions}

Each time you wish to return to the development environment, run the
following from the root of the repository:

\begin{lstlisting}[language=bash]
nix --extra-experimental-features 'nix-command flakes' develop
\end{lstlisting}

% ─────────────────────────────────────────────────────────────
\chapter{Database Setup}

\section{Automated Setup via Nix (Recommended)}

The recommended Nix workflow does not require a separately installed
PostgreSQL server. The \texttt{nix run .\#init}, \texttt{nix run .\#run},
\texttt{nix run .\#tests}, \texttt{nix run .\#seed}, and
\texttt{nix run .\#unseed} commands start a project-local PostgreSQL
cluster in \texttt{.postgres/data}, create the \texttt{flax10}
database, and export the required \texttt{POSTGRES\_*} variables
automatically.

\section{Manual PostgreSQL Setup}

If you are setting up the project without Nix, install and configure
PostgreSQL first.

\begin{enumerate}
    \item Create a PostgreSQL database and user matching the Django
          defaults in \texttt{settings.py}, or set the corresponding
          environment variables manually:
\begin{lstlisting}[language=sql]
CREATE USER postgres WITH PASSWORD '1234';
CREATE DATABASE postgres OWNER postgres;
\end{lstlisting}

    \item Alternatively, export custom connection variables before
          running Django commands:
\begin{lstlisting}[language=bash]
export POSTGRES_DB=postgres
export POSTGRES_USER=postgres
export POSTGRES_PASSWORD=1234
export POSTGRES_HOST=localhost
export POSTGRES_PORT=5432
\end{lstlisting}

    \item Apply migrations and seed the database from within
          \texttt{flax-10-project-files/backend}:
\begin{lstlisting}[language=bash]
python manage.py migrate
python manage.py seed
\end{lstlisting}
\end{enumerate}

\textbf{Note:} The project reads database connection values from
environment variables when provided. The default fallback values are
\texttt{postgres}/\texttt{postgres}/\texttt{1234}/
\texttt{localhost}/\texttt{5432}.

% ─────────────────────────────────────────────────────────────
\chapter{Manual Setup Without Nix}

This section is intended for developers who want to work without the
Nix flake.

\section{Backend}

\begin{enumerate}
    \item Create and activate a virtual environment from the repository
          root:
\begin{lstlisting}[language=bash]
python3.12 -m venv flax-10-project-files/venv
source flax-10-project-files/venv/bin/activate
\end{lstlisting}

    \item Install Python dependencies:
\begin{lstlisting}[language=bash]
pip install -r requirements.txt
\end{lstlisting}

    \item Configure PostgreSQL as described in Chapter~5.

    \item Run migrations and seed data:
\begin{lstlisting}[language=bash]
cd flax-10-project-files/backend
python manage.py migrate
python manage.py seed
\end{lstlisting}
\end{enumerate}

\section{Frontend}

\begin{enumerate}
    \item Install frontend dependencies:
\begin{lstlisting}[language=bash]
cd flax-10-project-files/frontend
npm install
\end{lstlisting}

    \item If required, provide frontend environment variables before
          starting Vite. The application can run without optional
          third-party API keys, but related map and travel-time
          features depend on them.
\end{enumerate}

% ─────────────────────────────────────────────────────────────
\chapter{Running the Application}

\section{Backend (Django)}

From within \texttt{flax-10-project-files/backend}:

\begin{lstlisting}[language=bash]
python manage.py runserver
\end{lstlisting}

The API will be available at \texttt{http://localhost:8000}.

\section{Frontend (React / Vite)}

From within \texttt{flax-10-project-files/frontend}:

\begin{lstlisting}[language=bash]
npm run dev
\end{lstlisting}

The frontend will be available at \texttt{http://localhost:5173}
(or the port shown in the Vite terminal output).

\textbf{Note:} When using the Nix workflow, start both servers with
\texttt{nix run .\#run}. That command selects the first available ports
starting from \texttt{8000} for the backend and \texttt{5173} for the
frontend.

% ─────────────────────────────────────────────────────────────
\chapter{Running the Tests}

\section{Backend Tests and Coverage}

From within \texttt{flax-10-project-files/backend}:

\begin{lstlisting}[language=bash]
python -m coverage erase
python -m coverage run --source=studyFlax manage.py test --verbosity=1
python -m coverage report -m
python -m coverage html
\end{lstlisting}

The HTML coverage report is generated at
\texttt{flax-10-project-files/backend/htmlcov/index.html} when run
manually.

\textbf{Note:} The Nix test entrypoint writes backend coverage to
\texttt{coverage/backend/html/index.html} and
\texttt{coverage/backend/coverage.xml}.

\section{Frontend Tests and Coverage}

From within \texttt{flax-10-project-files/frontend}:

\begin{lstlisting}[language=bash]
npm run test
npm run coverage
\end{lstlisting}

When using the Nix test entrypoint, frontend coverage is copied to
\texttt{coverage/frontend/}.

% ─────────────────────────────────────────────────────────────
\chapter{Environment Variables and Third-Party Services}

\section{Backend Variables}

\begin{center}
\begin{longtable}{lll}
\toprule
\textbf{Variable} & \textbf{Required} & \textbf{Purpose} \\
\midrule
POSTGRES\_DB & Yes & PostgreSQL database name \\
POSTGRES\_USER & Yes & PostgreSQL username \\
POSTGRES\_PASSWORD & Yes & PostgreSQL password \\
POSTGRES\_HOST & Yes & PostgreSQL host \\
POSTGRES\_PORT & Yes & PostgreSQL port \\
FRONTEND\_URL & Recommended & Used in password reset links \\
OPENAI\_API\_KEY & Optional & Enables AI analytics features \\
\bottomrule
\end{longtable}
\end{center}

\section{Frontend Variables}

\begin{center}
\begin{longtable}{lll}
\toprule
\textbf{Variable} & \textbf{Required} & \textbf{Purpose} \\
\midrule
VITE\_API\_BASE\_URL & Recommended & Backend API base URL \\
VITE\_MAPBOX\_TOKEN & Optional & Mapbox-powered map and geocoding features \\
VITE\_TFL\_API\_KEY & Optional & TfL travel-time integration \\
\bottomrule
\end{longtable}
\end{center}

\section{Third-Party Services}

\begin{itemize}
    \item \textbf{Mapbox} --- required for campus map and location
          autocomplete features. Provide a valid
          \texttt{VITE\_MAPBOX\_TOKEN}.
    \item \textbf{TfL API} --- required for travel-time features.
          Provide a valid \texttt{VITE\_TFL\_API\_KEY}.
    \item \textbf{OpenAI API} --- optional. AI analytics endpoints are
          only fully functional when \texttt{OPENAI\_API\_KEY} is set.
    \item \textbf{SMTP email} --- password reset emails use the SMTP
          configuration currently present in \texttt{settings.py}. For
          any real deployment, replace the default email account details
          with deployment-specific credentials.
\end{itemize}

% ─────────────────────────────────────────────────────────────
\chapter{User Fixtures and Access Credentials}

The following fixture accounts are available after seeding the
database. All passwords are \texttt{1234}.

\begin{center}
\begin{tabular}{lll}
\toprule
\textbf{Username} & \textbf{Role} & \textbf{Password} \\
\midrule
johndoe  & Student / Admin & 1234 \\
janedoe  & Student         & 1234 \\
amiller  & Teacher         & 1234 \\
charlie  & Student         & 1234 \\
\bottomrule
\end{tabular}
\end{center}

Additional users are seeded automatically using Faker.

\bigskip
\textbf{Django admin panel:}
\texttt{http://localhost:8000/admin/}

Signing in as \texttt{johndoe} with password \texttt{1234} also grants
access to admin features through the main application interface.

% ─────────────────────────────────────────────────────────────
\chapter{Deployed System}

The production deployment is hosted on Railway:

\begin{center}
\url{https://flax-10-production.up.railway.app/}
\end{center}

\section{Deployment Notes}

The repository includes a top-level \texttt{Dockerfile} that builds the
frontend and backend into a single container image. The container build
accepts the following build arguments for the frontend:

\begin{lstlisting}[language=bash]
VITE_MAPBOX_TOKEN
VITE_TFL_API_KEY
OPENAI_API_KEY
\end{lstlisting}

For deployment, ensure that the runtime environment also provides the
backend database connection variables
(\texttt{POSTGRES\_DB}, \texttt{POSTGRES\_USER},
\texttt{POSTGRES\_PASSWORD}, \texttt{POSTGRES\_HOST},
\texttt{POSTGRES\_PORT}) and \texttt{FRONTEND\_URL} if password reset
emails should point to the deployed frontend rather than localhost.

\textbf{Suggested deployment procedure for IT staff:}
\begin{enumerate}
    \item Provision a PostgreSQL database.
    \item Set the backend environment variables listed in Chapter~9.
    \item Provide the frontend build arguments required by the enabled
          third-party features.
    \item Build the Docker image using the repository \texttt{Dockerfile}.
    \item Run Django migrations:
\begin{lstlisting}[language=bash]
python manage.py migrate
\end{lstlisting}
    \item Optionally seed demonstration data:
\begin{lstlisting}[language=bash]
python manage.py seed
\end{lstlisting}
\end{enumerate}

% ─────────────────────────────────────────────────────────────
\chapter{Troubleshooting}

\begin{description}

    \item[Nix experimental features error.]
    If you encounter an error relating to \texttt{nix-command} or
    \texttt{flakes}, ensure the flags are passed as shown in
    Chapter~4, or add them permanently to \texttt{/etc/nix/nix.conf}.

    \item[Database connection error.]
    Verify that your PostgreSQL server is running and that the
    \texttt{POSTGRES\_*} environment variables match the target
    database. For non-Nix local setup, the fallback defaults are
    \texttt{postgres}/\texttt{postgres}/\texttt{1234}/
    \texttt{localhost}/\texttt{5432}.

    \item[Port already in use.]
    If \texttt{localhost:8000} or \texttt{localhost:5173} is already
    occupied, terminate the conflicting process or start the server
    on an alternative port:
\begin{lstlisting}[language=bash]
python manage.py runserver 8080
\end{lstlisting}

    For the frontend, you can use:
\begin{lstlisting}[language=bash]
npm run dev -- --port 5174
\end{lstlisting}

    \item[Coverage module not found.]
    Ensure \texttt{coverage} is installed in the virtual environment:
\begin{lstlisting}[language=bash]
pip install coverage
\end{lstlisting}

    \item[Re-seeding inflates data.]
    Running \texttt{python manage.py seed} repeatedly may duplicate seed
    data. To reset the database to a clean state, run the following
    from within \texttt{flax-10-project-files/backend}:
\begin{lstlisting}[language=bash]
python manage.py flush --no-input
python manage.py migrate
python manage.py seed
\end{lstlisting}

    If you are using the Nix workflow, \texttt{nix run .\#unseed} can
    also be used to remove seeded data before running
    \texttt{nix run .\#seed}.

\end{description}

% ─────────────────────────────────────────────────────────────
\chapter{Package Summary}

\section{Backend}

\begin{center}
\begin{longtable}{lll}
\toprule
\textbf{Package} & \textbf{Version} & \textbf{Purpose} \\
\midrule
Django              & 5.2.0  & Web framework \\
djangorestframework & 3.16.1 & REST API framework \\
psycopg             & 3.3.2  & PostgreSQL driver \\
psycopg-binary      & 3.3.2  & Precompiled psycopg binaries \\
faker               & 40.4.0 & Fake data generation for seeding \\
python-dotenv       & 1.2.1  & Environment variable management \\
openai              & 2.21.0 & OpenAI API integration \\
httpx               & 0.28.1 & HTTP client library \\
whitenoise          & 6.11.0 & Static file serving in deployment \\
\bottomrule
\end{longtable}
\end{center}

\section{Frontend}

\begin{center}
\begin{longtable}{lll}
\toprule
\textbf{Package} & \textbf{Version} & \textbf{Purpose} \\
\midrule
react           & 19.2.0  & UI library \\
react-dom       & 19.2.0  & React DOM renderer \\
vite            & 7.2.4   & Build tool and dev server \\
framer-motion   & 12.34.3 & Animation library \\
mapbox-gl       & 3.18.1  & Map rendering \\
styled-components & 6.3.9 & Component-scoped styling \\
eslint          & 9.39.1  & JavaScript linter \\
\bottomrule
\end{longtable}
\end{center}

\section{Package Management}

\begin{center}
\begin{tabular}{lll}
\toprule
\textbf{Tool} & \textbf{Version} & \textbf{Purpose} \\
\midrule
Nix & 2.33.2 & Environment package management \\
npm & 10.8.2 & JavaScript package management \\
\bottomrule
\end{tabular}
\end{center}

\end{document}
